% TODO: write abstract
Distances between probability measures are used across several areas of statistics, including the theory of estimation, hypothesis testing, and computational statistics. They provide natural tools for studying the convergence of estimators, quantifying dissimilarities between observed samples, and measuring discrepancies between empirical and target distributions. While classical applications mainly concern probability measures on Euclidean spaces, there is a need for distances on more complex spaces, such as spaces of laws of random probability measures. This is motivated by Bayesian nonparametrics, where the objects of inference are often random probability measures, and by modern statistical problems in which data objects themselves are naturally represented as probability measures. This thesis focuses on two applications in which such distances play a key role: two-sample testing and merging of opinions.

Two-sample testing assesses whether two populations differ by comparing observed samples, with the Kolmogorov–Smirnov test as a classic example. While numerous extensions address multivariate data, fewer methods have been developed for settings in which the objects of interest are complex, such as probability distributions themselves. In such a case, the random probability measures are typically not observed directly, but only through samples drawn from them. We propose a distance-based two-sample test for testing equality of laws of random probability measures using optimal transport theory, and leverage tools from empirical process theory to establish nonparametric theoretical guarantees. Empirically, we benchmark our method against existing approaches on simulated datasets and apply it to a mortality dataset.

The merging of opinions concerns whether two posterior distributions arising from different prior distributions eventually converge as more data are observed. Distances between posterior distributions provide a quantitative way to study both the occurrence of merging and its rate. We consider a Bayesian nonparametric framework, namely laws induced by completely random measures. Accordingly, the distance is tailored to completely random measures and is defined at the level of their Lévy intensities. Complementing existing results for other Bayesian nonparametric priors, this thesis derives merging rates between posterior distributions induced by the Pitman-Yor process and the Dirichlet process.
